% =============================================================================
%  2. Related work
%  Organised on OUR pipeline's axis (propose regions -> judge regions) rather
%  than by chronology, so that every cited work attaches to a claim we make.
%  Compressed from docs/paper/related_work.md, which keeps the full citation
%  ledger and the verbatim passages behind each statement.
%
%  NOTE FOR THE MDPI CONVERSION: MDPI papers in this family (e.g. Golej et al.,
%  Appl. Sci. 2024) fold related work into the Introduction and a "2.1
%  Background" subsection rather than giving it a numbered section of its own.
%  If that is the target, move 2.1-2.2 into the Introduction and 2.3-2.6 into
%  Materials and Methods / Background; the text below is written so that it
%  survives the move without rewriting.
% =============================================================================
\section{Related Work}
\label{sec:related}

\subsection{Anomaly detection without language}
\label{sec:related:novlm}

The industrial line of work defines the modern setting: train on nominal
images only, score anomalies at image and pixel level, and evaluate on
MVTec~AD~\cite{Bergmann_2019}. Its strongest representatives are memory-based
--- PatchCore~\cite{Roth_2022} stores a coreset of nominal patch features and
scores by nearest-neighbour distance --- distribution-based
--- PaDiM~\cite{Defard_2021} fits a per-position Gaussian over pretrained
features --- or distillation-based, where EfficientAD~\cite{Batzner_2024}
pushes latency to the millisecond range. The anomalib
library~\cite{Akcay_2022} makes these directly comparable, and we use its
reference implementations as baselines (\secref{sec:results:baselines}).
Two properties of this family matter here. First, it is object-centric: images
are crops of one category, not scenes. Second, it produces a score map, not a
region proposal with a semantic verdict, so it answers ``how anomalous'' rather
than ``what is this''.

Multi-class unified models remove the one-model-per-category assumption.
Dinomaly~\cite{Guo_2025} reconstructs frozen DINOv2~\cite{Oquab_2023} features
through a deliberately weakened decoder --- a noisy bottleneck, linear
attention, and a loose reconstruction constraint --- so that the model cannot
learn to reproduce anomalies it has never seen, and reports that a single
model matches per-category specialists. Its recent
successor~\cite{Guo_2025_dinomaly2} adds context-aware recentering of the
reconstruction targets among other changes. We re-implemented both and report
what they do on our cameras in \secref{sec:results:transfer}, including a
transfer result that this literature does not measure, because its benchmarks
have no notion of a \emph{viewpoint} that a model could fail to transfer
across.

Closest to our proposal stage is AnomalyDINO~\cite{Damm_2025}, which shows that
frozen DINOv2 patch features with a nearest-neighbour memory bank are a
strong training-free anomaly detector and, notably, that this vision-only
approach outperforms vision--language alternatives in the few-shot regime. Our
proposal stage adopts its scoring idea and specialises it to a fixed camera
(\secref{sec:method:dino}).

\subsection{Change detection and abandoned objects}
\label{sec:related:aod}

Comparing a frame against a nominal reference is change detection, and CDnet
2014~\cite{Wang_2014} is its reference benchmark: eleven categories of
perturbation --- shadows, camera jitter, night, dynamic background,
intermittent object motion --- with per-pixel ground truth and a fixed metric
suite. Its \texttt{intermittentObjectMotion} category is our task almost word
for word, and we use CDnet to evaluate the proposal stage in isolation
(\secref{sec:results:cdnet}).

Abandoned-object detection (AOD) is the task whose \emph{semantics} match ours.
Luna et al.~\cite{Luna_2018} survey the field and propose a unified evaluation
protocol; the representative modern pipeline still combines dual background
models with a CNN classifier~\cite{Park_2020,Park_2019}, and recent
benchmarks~\cite{Xu_2024} stay outdoors and road-side. Two things stand out for
our purposes. The architecture is unchanged since the AVSS
era~\cite{Miguel_2008} --- background subtraction proposes, a closed-vocabulary
classifier decides --- and illumination change remains the openly stated
failure mode. Our pipeline replaces the classifier with an open-vocabulary
judge and the illumination problem with an explicit occlusion mask plus a
feature-space comparison that is largely exposure-invariant.

\subsection{Language-guided anomaly detection}
\label{sec:related:lvlm}

One family puts language inside the scoring function.
WinCLIP~\cite{Jeong_2023} scores windows against handcrafted normal/anomalous
prompts; AnomalyCLIP~\cite{Zhou_2023} learns object-agnostic prompt vectors so
that one prompt set transfers across datasets. A second family uses a large
vision--language model as the detector itself: AnomalyGPT~\cite{Gu_2024}
produces threshold-free textual verdicts and reports that such models are weak
on \emph{localized} detail --- which is precisely why our judge is never asked
to localize. The survey of Ren et al.~\cite{Ren_2025} organises the field into
foundation models used as encoder, as detector, or as interpreter, and names
prompt sensitivity as a recurring obstacle. Our judge is an
\emph{interpreter} in that taxonomy: it never produces the box.

\subsection{Video anomaly detection and training-free VLM pipelines}
\label{sec:related:vad}

The weakly supervised mainstream, anchored by UCF-Crime~\cite{Sultani_2018} and
surveyed in~\cite{Duong_2023,Li_2025}, learns from video-level labels and
outputs a per-frame anomaly score. It is not our benchmark: its anomalies are
events with temporal extent, ours are static changes that persist
indefinitely, and its cameras are uncontrolled.

Three recent systems are our closest competitors because they, like us, use a
vision--language model without task-specific training.
LAVAD~\cite{Zanella_2024} captions frames with an off-the-shelf captioner,
cleans the captions with an LLM, and scores temporally.
VERA~\cite{Ye_2025} learns \emph{verbalized} guiding questions as parameters
and aggregates segment scores into frame scores.
AnomalyRuler~\cite{Yang_2024} induces normality rules from a few normal frames
and applies them as language. All three decide \emph{when} something is
anomalous, and all three consume \emph{whole frames}. AnomalyRuler is explicit
about it: its visual perception module ``utilizes a VLM to convert video frames
into text descriptions'', prompted with \emph{``What are people doing? What are
in the images other than people?''} --- a question about the frame, not about a
region of it. None of the three contains a spatial region-proposal stage that
could be varied while the language model is held fixed, which is why the
experiment of \secref{sec:results:main} has no counterpart in this line of
work.

\subsection{Can the judge be trusted?}
\label{sec:related:judge}

Our judge answers a binary question about a crop, which is exactly the probing
format of POPE~\cite{Li_2023}: POPE shows that state-of-the-art vision--language
models hallucinate objects at high rates under yes/no probing, and that the
measured rate depends strongly on how the negative examples are sampled. That
is the reason we report every recall next to the judge's YES rate
(\secref{sec:setup:metrics}): a model that answers YES to most crops recovers
most instances mechanically, and calling that detection would be an artefact of
the sampling. The broader concern of using a model as an evaluator is by now a
named methodological problem~\cite{Gu_2026}.

\subsection{Region proposal and in-vehicle CCTV}
\label{sec:related:proposal}

Open-vocabulary detectors --- Grounding DINO~\cite{Liu_2023},
YOLO-World~\cite{Cheng_2024} --- propose boxes from text and are an obvious
alternative proposal stage; SAHI~\cite{Akyon_2022} shows that tiled inference
substantially improves small-object detection, which is the regime a forgotten
phone falls into at CCTV resolution.

Closest to our optional veto stage is BEM~\cite{Park_2026_bem}, a
training-free, weight-frozen module attached to a pretrained detector at
inference: it estimates a background embedding from recent unlabelled frames of
a fixed camera, keeps a prototype memory, and re-scores the detector's logits
with an inverse-similarity penalty, reducing false positives while preserving
recall. Its premise is ours --- on a fixed camera the quasi-static background is
a stable, label-free prior --- and its position in a pipeline is ours too: it
acts \emph{after} a proposal stage and can only suppress. It is therefore
independent corroboration of one half of our finding
(\secref{sec:results:main}), from a different task and a different detector
family. Two things separate it from our work: its prior is a single global
frame embedding rather than a per-position reference, and it re-scores the
outputs of a closed-vocabulary detector, so what it can never do is put a box
around something the detector did not already propose.

The in-vehicle deployment domain itself is thinly covered, and the review by
Caetano et al.~\cite{Caetano_2022} says so in terms that match our setting
closely. It identifies the extension of outdoor abandoned-object and
abnormal-behaviour work to in-vehicle monitoring as ``a relevant research
opportunity that has been overlooked in the accessible literature'', names
moving backgrounds and frequent illumination changes as the specific issues
previous surveys ignore, and reports that work in this domain remains ``fully
dependent on the availability of private datasets''. It also states the
operational priority --- the identification of unattended objects as the main
safety challenge in public transport --- and the very heuristic our person veto
implements: an object counts as abandoned when its presence is not expected
\emph{without a person present in the frame}. Its proposed remedy for the data
scarcity is synthetic augmentation, which is the family our inpainted cases
belong to (\secref{sec:setup:data}). On the benchmark side, Bus
Violence~\cite{Ciampi_2022} is, to our knowledge, the only public in-vehicle
CCTV anomaly benchmark, and it targets violent \emph{behaviour} --- an event,
with people as the subject --- rather than persistent change with people as the
principal nuisance. The gap our data occupies is therefore documented, not
assumed, and it is also the reason we cannot fill it publicly
(\secref{sec:discussion}).
